\documentclass[10pt,a4paper]{amsart}
\usepackage[margin=19mm]{geometry}
\usepackage{amsmath,amssymb,mathtools}
\usepackage[scr=rsfs,cal=euler]{mathalfa}
\usepackage{xfrac}
\usepackage[unicode,hidelinks,bookmarks=false]{hyperref}
\newcommand{\ie}{i.e.\ }
\newcommand{\cf}{cf.\ }
\newcommand{\resp}{respectively\ }
\newcommand{\Subsection}{\subsection}
\providecommand{\nobreakdash}{\nobreak\hbox{-}}
\setlength{\emergencystretch}{2em}
\multlinegap0pt
\usepackage{amssymb}
\usepackage{mathtools}
\usepackage[shortlabels]{enumitem}
\newtheorem{theorem}{Theorem}
\title{Non-colliding billiards in the plane}
\author{Itai Benjamini, Alexander Shamov, Barak Weiss}
\date{Source: arXiv:2605.23575v2 (CC BY 4.0)}
\begin{document}
\maketitle

\noindent\emph{Source and licence.} Non-colliding billiards in the plane. Version arXiv:2605.23575v2 (CC BY 4.0), distributed by arXiv under the Creative Commons Attribution 4.0 International licence, http://creativecommons.org/licenses/by/4.0/, as recorded in the arXiv licence metadata for this exact version and shown on the abstract page https://arxiv.org/abs/2605.23575v2. Full licence notice: \texttt{licenses/arxiv-2605.23575-CC-BY-4.0.txt}.\par\smallskip

\noindent\textit{Editorial scope.} Counted unit: the article's theorem thm:no-universal (source0.tex lines 248--254). The article's complete original proof of it is delivered with it (source0.tex lines 256--422, one \texttt{proof} environment of 167 lines). No mathematics is written, repaired or re-derived: the statement and the proof are the article's own lines, in the article's own notation, and no retained line is substituted or re-flowed.  No further source-local context is needed: every cross-reference of every retained line resolves inside the two delivered spans.  Nothing was omitted from any retained span, so the package carries no omission record.  No retained line contains a citation, so the wrapper carries no bibliography and no bibliography entry.  No retained line contains a figure environment, an image-inclusion command or drawing code, so no image is delivered and the package declares no asset dependency.  Editorial formatting only, in addition: an AMS \texttt{amsart} preamble that restates the article's own declarations and macros used by the retained lines, namely the environments \texttt{theorem} on separate counters, together with the standard packages those lines need.  The retained environments print in this document as Theorem 1 in order of appearance, whereas the article prints the same items with the numbering of its own full document.  The complete native source of the article as distributed in the arXiv e-print for this exact version is bundled unchanged as \texttt{sources/201221/source0.tex}.
\begin{theorem}\label{thm:no-universal}
There is no bounded continuous function $w:\mathbb R^2\to\mathbb R^2$ for which there exists a constant $c>0$ such that
\begin{equation}\label{eq:universal-abs}
\bigl|\langle x-y,w(x)-w(y)\rangle\bigr|>c\|w(x)-w(y)\|
\end{equation}
for every pair $x,y\in\mathbb R^2$ with $\|x-y\|>1$.
\end{theorem}

\begin{proof}
Assume, for a contradiction, that such a function $w$ exists.

Let
\[
D:=\{(x,y)\in\mathbb R^2\times\mathbb R^2:\|x-y\|>1\}.
\]
This set is connected. Indeed, under the change of variables $(x,y)\mapsto ((x+y)/2,x-y)$, it is the product of $\mathbb R^2$ with the connected set $\{z\in\mathbb R^2:\|z\|>1\}$.

The inequality \eqref{eq:universal-abs} implies that
\[
\langle x-y,w(x)-w(y)\rangle\ne 0
\]
for every $(x,y)\in D$. Since this inner product is continuous on the connected set $D$, its sign is constant. Replacing $w$ by $-w$ if necessary, we may assume that
\begin{equation}\label{eq:universal-positive}
\langle x-y,w(x)-w(y)\rangle>c\|w(x)-w(y)\|
\end{equation}
for all $(x,y)\in D$.

Set
\[
a:=\min\{c/2,1/2\}\in(0,1),
\qquad
\delta:=\arcsin a>0.
\]
For every nonzero vector $u\in\mathbb R^2$, define the closed cone
\[
C_u:=\{z\in\mathbb R^2:\langle u,z\rangle\ge a\|u\|\,\|z\|\}.
\]
Equivalently, $C_u$ consists of all vectors whose angle with $u$ is at most
\[
\arccos a=\frac{\pi}{2}-\delta .
\]
In particular, each $C_u$ is a closed convex cone, and $C_{\lambda u}=C_u$ for every $\lambda>0$.

We first record a consequence of \eqref{eq:universal-positive}. If $1<\|x-y\|\le 2$, then
\[
c\ge a\|x-y\|,
\]
and therefore \eqref{eq:universal-positive} implies
\begin{equation}\label{eq:cone-near}
w(x)-w(y)\in C_{x-y}.
\end{equation}
We claim that the same conclusion holds for all $x,y$ with $\|x-y\|>1$:
\begin{equation}\label{eq:cone-all}
w(x)-w(y)\in C_{x-y}.
\end{equation}
Let $L=\|x-y\|>1$. Choose an integer $n$ such that
\[
L/2\le n<L;
\]
such an integer exists, for example $n=\lceil L/2\rceil$. Then $L/n\in(1,2]$. Put
\[
z_i=x+\frac{i}{n}(y-x), \qquad i=0,1,\ldots,n.
\]
For each $i$, the segment length $\|z_i-z_{i+1}\|=L/n$ lies in $(1,2]$, and $z_i-z_{i+1}$ is a positive scalar multiple of $x-y$. Hence \eqref{eq:cone-near} gives
\[
w(z_i)-w(z_{i+1})\in C_{x-y}.
\]
Since $C_{x-y}$ is a convex cone,
\[
w(x)-w(y)=\sum_{i=0}^{n-1}\bigl(w(z_i)-w(z_{i+1})\bigr)
\in C_{x-y},
\]
proving \eqref{eq:cone-all}.

Now consider the cluster set of $w$ at infinity,
\[
V_\infty:=\bigcap_{R>0}
\overline{\,w\bigl(\{x\in\mathbb R^2:\|x\|\ge R\}\bigr)\,}.
\]
Here and below, the bar denotes closure. Because $w$ is bounded, each set in the intersection is compact. It is also connected: the exterior $\{x:\|x\|\ge R\}$ is connected, its image under the continuous map $w$ is connected, and the closure of a connected set is connected. These compact connected sets are nested as $R$ increases. Therefore $V_\infty$ is nonempty, compact, and connected.

We shall use the following elementary characterization: a point $\xi\in\mathbb R^2$ lies in $V_\infty$ if and only if there is a sequence $x_n\in\mathbb R^2$ such that $\|x_n\|\to\infty$ and $w(x_n)\to\xi$. Indeed, one implication is immediate from the definition. Conversely, if $\xi\in V_\infty$, then for each $n$ we may choose $x_n$ with $\|x_n\|\ge n$ and $\|w(x_n)-\xi\|<1/n$.

The set $V_\infty$ is not a singleton. Suppose instead that $V_\infty=\{\xi\}$. Then $w(x)\to\xi$ as $\|x\|\to\infty$; otherwise, boundedness of $w$ would produce a sequence tending to infinity along which $w$ has a subsequential limit different from $\xi$, contradicting the definition of $V_\infty$.

Fix $y\in\mathbb R^2$ and a unit vector $u\in S^1$. Applying \eqref{eq:cone-all} to $x=ru$ and this fixed $y$, and then letting $r\to\infty$, gives
\[
\xi-w(y)\in C_u.
\]
Since this holds for every $u\in S^1$, it follows that $\xi-w(y)=0$; for example, the conditions for $u$ and $-u$ can both hold only for the zero vector. Thus $w(y)=\xi$ for every $y$, so $w$ is constant. This contradicts the strict inequality \eqref{eq:universal-abs}. Hence $V_\infty$ is a nontrivial connected compact set. In particular, it contains arbitrarily many distinct points.

It remains to obtain a contradiction from the angular restrictions imposed by \eqref{eq:cone-all}. Let $\mathbb T=\mathbb R/2\pi\mathbb Z$, and let $d_{\mathbb T}$ denote circular distance on $\mathbb T$. For a nonzero vector $z$, write $\arg z\in\mathbb T$ for its direction.

We claim that if $\xi,\eta\in V_\infty$ are distinct, and if there are sequences $x_n,y_n$ with
\[
\|x_n\|\to\infty,
\qquad
\|y_n\|\to\infty,
\qquad
w(x_n)\to\xi,
\qquad
w(y_n)\to\eta,
\]
such that
\[
\arg x_n\to\alpha,
\qquad
\arg y_n\to\beta
\]
for some $\alpha,\beta\in\mathbb T$, then
\begin{equation}\label{eq:angular-separation}
d_{\mathbb T}(\alpha,\beta)\ge 2\delta .
\end{equation}

To prove the claim, choose indices $N(n)$ tending to infinity so fast that
\[
\frac{\min\{\|x_{N(n)}\|,\|y_{N(n)}\|\}}
{\max\{\|x_n\|,\|y_n\|\}}
\longrightarrow \infty .
\]
Then
\[
\arg (y_{N(n)}-x_n)\to\beta,
\qquad
\arg (x_{N(n)}-y_n)\to\alpha .
\]
For all large $n$, the relevant distances are greater than $1$, so \eqref{eq:cone-all} gives
\[
w(y_{N(n)})-w(x_n)\in C_{y_{N(n)}-x_n},
\qquad
w(x_{N(n)})-w(y_n)\in C_{x_{N(n)}-y_n}.
\]
Let $q:=\eta-\xi\ne 0$. Since
\[
w(y_{N(n)})-w(x_n)\to q,
\qquad
w(x_{N(n)})-w(y_n)\to -q,
\]
the preceding cone inclusions imply
\[
d_{\mathbb T}(\arg (y_{N(n)}-x_n),\arg q)
\le \frac{\pi}{2}-\delta+o(1),
\]
and
\[
d_{\mathbb T}(\arg (x_{N(n)}-y_n),\arg(-q))
\le \frac{\pi}{2}-\delta+o(1).
\]
Because $d_{\mathbb T}(\arg q,\arg(-q))=\pi$, the triangle inequality on the circle yields
\[
d_{\mathbb T}(\arg (y_{N(n)}-x_n),\arg (x_{N(n)}-y_n))
\ge 2\delta-o(1).
\]
Letting $n\to\infty$ proves \eqref{eq:angular-separation}.

Now choose an integer $k>\pi/\delta$ and choose distinct points
\[
\xi_1,\ldots,\xi_k\in V_\infty .
\]
For each $i$, by the sequence characterization of $V_\infty$ and compactness of the circle, choose a sequence $x_n^{(i)}$ such that
\[
\|x_n^{(i)}\|\to\infty,
\qquad
w(x_n^{(i)})\to\xi_i,
\qquad
\arg x_n^{(i)}\to\alpha_i\in\mathbb T .
\]
The claim gives
\[
d_{\mathbb T}(\alpha_i,\alpha_j)\ge 2\delta
\]
for all $i\ne j$. But the circle of circumference $2\pi$ cannot contain more than $\pi/\delta$ points whose pairwise circular distances are all at least $2\delta$: the open arcs of radius $\delta$ around such points are pairwise disjoint. This contradicts $k>\pi/\delta$.

The contradiction proves the theorem.
\end{proof}

\end{document}
